\begin{textsample}{NEG Dim 4 – 2016\_c5 – Score 2.00 – 2016\_c5/enade\_2016\_agro\_6f2e9fbda0.txt}  \label{ex:f4_neg_254}
CAPÍTULO 5 – DISTRIBUIÇÃO DOS CONCEITOS 5.1 - PANORAMA NACIONAL DA DISTRIBUIÇÃO DOS CONCEITOS O Conceito Enade é um indicador de qualidade que avalia o desempenho dos estudantes a partir dos resultados do Enade.
Os valores possíveis vão de 1 ( pior situação ) a 5 ( melhor situação ).
As Unidades de Observação com apenas um ou sem nenhum concluinte participante não obtêm o Conceito Enade, ficando Sem Conceito ( SC ).
A Tabela 5.1 apresenta a quantidade e distribuição de cursos de Agronomia participantes do Enade/2016, por faixa de conceito e Grande Região.
Enfatiza -se, mais uma vez, que a diferença entre os cursos tabulados neste capítulo e no capítulo 2 corresponde aos cursos sem conceito, ou seja àqueles, em princípio, sem alunos concluintes que houvessem participado da prova ou àqueles que tivessem tão somente um aluno concluinte19.
Observando -se os dados da Tabela 5.1, nota -se que, dos 253 cursos participantes, 88 ( 34,8 \% ) classificaram -se com conceito 3, e outros 88 com conceito 4, os valores modais.
O conceito 3 foi o conceito modal em três regiões : Norte ( 50,0 \% ), Nordeste ( 31,7 \% ) e Centro-Oeste ( 50,0 \% ).
Nas demais regiões, a moda foi o conceito 4, atribuído para 44,9 \% dos cursos da região Sudeste e para 46,8 \% dos cursos da região Sul.
O conceito 2 foi o segundo mais frequente em nível nacional ( 19,4 \%, correspondendo a 49 cursos ) e em seguida os conceitos 1 e 5, ( 5,5 \%, correspondendo a 14 cursos cada ).
Dos 253 cursos de Agronomia, nenhum ( 0,0 \% ) ficou sem conceito ( SC ). 19 Estes cursos com somente um concluinte não foram considerados no capítulo 2.
A região Norte participou com 28 cursos ou 11,1 \% do total nacional.
Desses, 14 cursos receberam o conceito 3, o que equivale a 50,0 \% do total regional, além de ser o conceito modal na região, como já comentado.
A seis cursos ( 21,4 \% ) atribuiu -se conceito 1, a cinco cursos ( 17,9 \% ), conceito 2, a dois cursos ( 7,1 \% ), conceito 4, e a um curso ( 3,6 \% ), conceito 5.
Nenhum curso ficou sem conceito ( SC ) nessa região.
A região Nordeste participou com 41 cursos ou 16,2 \% do total nacional.
Nessa região, 13 cursos ( 31,7 \% em termos regionais ) obtiveram conceito 3, o conceito modal para a região, como já comentado.
Os conceitos 1 e 2 foram atribuídos, respectivamente, a quatro e a 12 cursos ( 9,8 \% e 29,3 \%, respectivamente ).
Ainda nessa região, 12 cursos ( 29,3 \% ) foram avaliados com o conceito 4, e nenhum curso ( 0,0 \% ), com conceito 5.
Nenhum curso ficou sem conceito ( SC ) nesta região tampouco.
Dos 78 cursos participantes da região Sudeste ( 30,8 \% do total nacional ), o conceito 4 foi o conceito modal, como já anteriormente assinalado, tendo sido obtido por 35 cursos ( 44,9 \% ).
O conceito 1 foi atribuído a dois cursos ( 2,6 \% ) e o conceito 2, a 18 cursos ( 23,1 \% ).
Receberam os conceitos 3 e 5, respectivamente, 21 cursos ( 26,9 \% ) e dois cursos ( 2,6 \% ).
Nessa região, tampouco, nenhum curso ficou sem conceito ( SC ).
Os 62 cursos da região Sul corresponderam a 24,5 \% do total nacional.
Como já indicado, houve predominância do conceito 4, atribuído a 46,8 \% dos cursos da região ( conceito modal ), ou seja, atribuído a 29 dos 62 cursos participantes na região Sul.
O conceito 2 foi atribuído a quatro cursos ( 6,5 \% ), o conceito 3, a 18 cursos ( 29,0 \% ) e o conceito 5, a 11 cursos ( 17,7 \% ).
Nenhum curso recebeu conceito 1 ou ficou sem conceito ( SC ) na região Sul.
Já dos 44 cursos participantes na região Centro-Oeste ( 17,4 \% do total nacional ), 22 receberam conceito 3, o conceito modal ( 50,0 \% ), como já fora destacado.
Por sua vez, o conceito 1 foi atribuído a dois cursos ( 4,5 \% ).
Já os conceitos 2 e 4 tiveram o mesmo número de cursos, dez ( 22,7 \% ), enquanto que o conceito 5 não foi atribuído a nenhum curso.
Nessa região, nenhum curso ficou sem conceito ( SC ).
O Gráfico 5.1 apresenta a distribuição do Conceito Enade segundo Grande Região : poligonais mais à esquerda representam Grandes Regiões com pior distribuição, e poligonais mais à direita, Grandes Regiões com melhores conceitos.
A região Norte ( linha azul ) apresenta os cursos com conceitos mais baixos, ao passo que a região Sul ( linha roxa ) apresenta a distribuição com valores maiores.
Cursos em IES situadas nas regiões Nordeste ( linha verde ) e Centro-Oeste ( linha preta ) apresentam uma situação intermediária.
A região Sudeste ( linha vermelha ) apresenta uma distribuição com conceitos melhores do que as regiões Norte, Nordeste e Centro-Oeste, mas pior do que a Sul. 5.2 - CONCEITOS POR CATEGORIA ADMINISTRATIVA E GRANDE REGIÃO A Tabela 5.2 apresenta a distribuição dos cursos participantes do Enade/2016 de Agronomia, por Categoria Administrativa, de acordo com os conceitos por eles alcançados, segundo as Grandes Regiões brasileiras.
Dos 253 cursos participantes, 152 ( 60,1 \% ) eram ministrados em Instituições Públicas e 101 ( 39,9 \% ), em Instituições Privadas.
De acordo com as informações da Tabela 5.2, em termos nacionais, entre Instituições Públicas e Privadas, observa -se que, dos 14 cursos avaliados com conceito 5, 13 eram oferecidos em IES Públicas e um, em IES Privadas.
Dos 152 cursos participantes de IES Públicas, o conceito 4 foi o valor modal, atribuído a 60 cursos ( 39,5 \% ).
Entre os demais cursos de Instituições Públicas participantes, seis cursos ( 3,9 \% ) foram avaliados com conceito 1, 20 cursos ( 13,2 \% da categoria ) obtiveram conceito 2, 53 cursos ( 34,9 \% ) receberam conceito 3 e, como já comentado, 13 cursos ( 8,6 \% ) receberam conceito 5.
Nesta categoria, nenhum curso ficou sem conceito ( SC ).
Na rede privada, o conceito modal foi 3, atribuído a 35 cursos, o correspondente a 34,7 \% dos 101 cursos da categoria.
Entre os demais cursos participantes, oito ( 7,9 \% ) receberam conceito 1, e 29 ( 28,7 \% ), conceito 2.
O conceito 4 foi atribuído a 28 cursos ( 27,7 \% ) e o conceito 5, a um curso ( 1,0 \% ).
Nesta Categoria Administrativa, nenhum curso ficou sem conceito ( SC ).
Na análise por região, observa -se que, na região Norte, as Instituições Públicas participaram com 21 cursos ( 75,0 \% do total regional ), dos quais nenhum ficou sem conceito ( SC ) e um ficou com conceito 5 ( 4,8 \% ).
O conceito modal para as Instituições Públicas na região foi 3, atribuído a 12 cursos, correspondendo a 57,1 \% dessa categoria na região.
Dois cursos foram avaliados com conceito 1, quatro cursos com conceito 2, e, finalmente, dois com conceito 4.
As Instituições Públicas participaram com sete cursos na região Norte ( 25,0 \% do total regional ), dos quais quatro obtiveram conceito 1, o valor modal ( 57,1 \% ).
O conceito 2 foi atribuído a um curso ( 14,3 \% ), enquanto o conceito 3 foi atribuído a dois cursos ( 28,6 \% ).
Já nenhum curso foi atribuído aos conceitos restantes, conceitos 4 e 5.
Além disso, nenhum curso ficou sem conceito ( SC ) nessa combinação de categoria e região.
Na região Nordeste, a rede pública concentrou 37 dos 41 cursos participantes, o equivalente a 90,2 \% do total da região, a maior das proporções desta categoria entre as regiões.
Dentre os cursos oferecidos por IES Públicas no Nordeste, 13 cursos foram avaliados com conceito 3, o conceito modal.
Um curso obteve conceito 1, 12 cursos receberam conceito 2, 11 cursos receberam conceito 4 e, finalmente, nenhum curso alcançou o conceito máximo, 5.
Além disso, nessa combinação de categoria e região, nenhum curso ficou sem conceito ( SC ).
As Instituições Privadas da região Nordeste participaram com quatro cursos ( 9,8 \% do total da região ), a menor das proporções desta categoria nas regiões.
Desses, 3 obtiveram conceito 1, o valor modal ( 75,0 \% ).
O outro curso recebeu conceito 4.
Os conceitos 2, 3 e 5 não foram atribuídos a nenhum curso.
Por fim, nenhum curso ficou sem conceito ( SC ) nessa região.
Na região Sudeste, a proporção de cursos da rede privada, 53,8 \%, foi mais elevada do que a registrada nas demais regiões brasileiras, correspondendo a 42 dos 78 cursos participantes.
Nessa categoria, o conceito modal foi 2, atribuído a 16 cursos ( 38,1 \% ), diferentemente do conceito modal da região que, como já mencionado, foi o conceito 4.
Além disso, nenhum curso ficou sem conceito ( SC ) nem recebeu conceito 1 ou 5.
Os demais foram avaliados com conceito 3 ( 15 cursos e 35,7 \% ) e com conceito 4 ( 11 cursos e 26,2 \% ).
Entre os 36 cursos oferecidos por Instituições Públicas na região Sudeste ( 46,2 \% do total regional ), o conceito modal foi 4, atribuído a 24 cursos ( 66,7 \% ).
Os demais cursos receberam os conceitos 1 ( dois cursos ), 2 ( dois cursos ), 3 ( seis cursos ) e, finalmente, 5 ( dois cursos ).
Na região Sudeste, nesta categoria, nenhum curso ficou sem conceito ( SC ).
As Instituições Públicas concentraram 31 dos 62 cursos participantes da região Sul, exatamente metade do total regional.
Desses, 15 obtiveram conceito 4, o conceito modal ( 48,4 \% ).
Nesta combinação de Categoria Administrativa e Grande Região, os conceitos 1 e 2 não foram atribuídos a nenhum curso.
Seis cursos receberam conceito 3, dez cursos receberam conceito 5 e nenhum ficou sem conceito ( SC ).
As Instituições Privadas na região Sul participaram com outros 31 cursos, a outra metade dos cursos da região.
O conceito modal foi 4, atribuído a 14 cursos.
Aos demais cursos foram atribuídos os conceitos 2 ( quatro cursos ), 3 ( 12 cursos ) e 5 ( um curso ).
Nenhum curso da região Sul oferecido por IES Privada ficou sem conceito ( SC ) ou recebeu conceito 1.
Na região Centro-Oeste, 27 dos 44 cursos participantes eram de Instituições Públicas ( 61,4 \% em termos regionais ).
Desses, 16 receberam conceito 3, o conceito modal.
Entre os demais, um recebeu conceito 1, dois receberam conceito 2, oito receberam conceito 4 e nenhum recebeu conceito 5.
Nenhum curso ficou sem conceito ( SC ).
Quanto aos 17 cursos oferecidos por Instituições Privadas na região Centro-Oeste ( 38,6 \% do total regional ), notam -se as maiores frequências em dois conceitos contíguos, os conceitos 2 ( conceito modal ) e 3, atribuídos, respectivamente, a oito e a seis cursos.
Um curso recebeu conceito 1, e dois cursos ficaram com conceito 4.
Nenhum curso ficou sem conceito ( SC ) ou recebeu conceito 5.
O Gráfico 5.2 apresenta a distribuição do Conceito Enade segundo a Categoria Administrativa da IES.
Os cursos em IES Públicas ( linha azul ) apresentam uma poligonal mais à direita do que os em IES Privadas ( linha verde ) e, consequentemente, uma distribuição de Conceitos Enade com valores maiores. 5.3 - CONCEITOS POR ORGANIZAÇÃO ACADÊMICA E GRANDE REGIÃO Na Tabela 5.3, encontra -se a distribuição dos conceitos atribuídos aos cursos participantes do Enade/2016 na Área de Agronomia, por Organização Acadêmica, segundo as Grandes Regiões brasileiras.
Dos 253 cursos de Agronomia participantes, 176 eram oferecidos em Universidades, 26 em Centros Universitários e 51 em Faculdades.
Esta distribuição corresponde a, respectivamente, 69,6 \%, 10,3 \% e 20,2 \% dos cursos.
De acordo com os dados apresentados, de todos os 14 cursos avaliados com conceito 5, 13 eram vinculados a Universidades.
Os cursos vinculados a esse tipo de Organização Acadêmica tiveram o conceito 4 como conceito modal, atribuído a 70 cursos ( 39,8 \% ).
Os demais cursos vinculados a Universidades receberam os conceitos 1 ( seis cursos ), 2 ( 21 cursos ), 3 ( 66 cursos ) e 5 ( 13 cursos, como já mencionado ).
Nenhum curso ficou sem conceito ( SC ) neste tipo de Organização Acadêmica.
Entre os cursos em Centros Universitários, o conceito modal foi 3, atribuído a nove cursos ( 34,6 \% ).
Os outros cursos receberam os conceitos 1 ( um curso ), 2 ( oito cursos ), 4 ( sete cursos ) e 5 ( um curso ).
Nenhum curso vinculado a esse tipo de Organização Acadêmica ficou sem conceito ( SC ).
Dos 51 cursos mantidos por Faculdades, 20 ( 39,2 \% ) receberam conceito 2, o conceito modal.
O segundo conceito mais frequentemente atribuído aos cursos oferecidos por essa categoria de IES foi o conceito 3, recebido por 13 cursos.
Na sequência, 11 cursos obtiveram conceito 4 e sete cursos, conceito 1.
Nenhum curso ficou sem conceito ( SC ) ou recebeu conceito 5.
Considerando -se separadamente as regiões brasileiras, verifica -se que, na região Norte, as Universidades concentraram 21 dos 28 cursos participantes.
Doze cursos receberam conceito 3, conceito modal.
Os demais cursos receberam os conceitos 1 ( dois cursos ), 2 ( quatro cursos ), 4 ( dois cursos ) e 5 ( um curso ).
Nenhum curso ficou sem conceito ( SC ) nesta combinação de categoria e região.
Os Centros Universitários da região Norte foram representados por dois cursos, aos quais foram atribuídos os conceitos 1 e 3.
Nenhum curso ficou sem conceito ( SC ) ou recebeu os demais conceitos nesta combinação de categoria e região.
As Faculdades participaram com 5 cursos na região Norte que receberam os conceitos 1 ( três cursos, conceito modal ), 2 ( um curso ) e 3 ( um curso ).
Nenhuma Faculdade da região ficou sem conceito ( SC ) ou obteve conceitos 4 e 5.
Na região Nordeste, as Universidades participaram com 36 dos 41 cursos da Área de Agronomia oferecidos na região, nenhum sem conceito ( SC ).
O conceito modal foi 3, atribuído a 13 cursos, enquanto os conceitos 2 e 4 foram atribuídos a 11 cursos cada.
O conceito 1 foi obtido por apenas um curso.
Nenhum curso oferecido por Universidades no Nordeste recebeu conceito 5.
Os Centros Universitários não contaram com cursos participantes na região Nordeste.
Já as Faculdades foram representadas por 5 cursos na região, 3 dos quais obtiveram o conceito modal 1.
O conceito 2 foi atribuído a um curso, assim como o conceito 4.
Na região Nordeste, nenhum curso oferecido por Faculdades recebeu conceitos 3 ou 5, e tampouco nenhum curso ficou sem conceito ( SC ).
Na região Sudeste, as Universidades concentraram 43 dos 78 cursos de Agronomia da região.
Entre os cursos em Universidades, nessa região, o conceito modal foi 4, atribuído a 27 cursos.
Os demais cursos receberam os conceitos 1 ( dois cursos ), 2 ( dois cursos ), 3 ( dez cursos ) e 5 ( dois cursos ).
Nenhum curso ficou sem conceito.
Os Centros Universitários participaram com 14 cursos na região Sudeste, dos quais seis obtiveram o conceito modal 2.
O conceito 3 foi atribuído a três cursos, o conceito 2, a cinco cursos, enquanto que os conceitos 1 e 5 não foram atribuídos a nenhum curso.
Além disso, nenhum curso ficou sem conceito ( SC ) nessa combinação de organização acadêmica e região.
Já as Faculdades foram representadas por 21 cursos na região Sudeste, sendo 10 atribuídos ao conceito modal 2.
Aos demais cursos foram atribuídos os conceitos 3 ( oito cursos ) e 4 ( três cursos ).
Nenhum curso ficou sem conceito ( SC ) ou foi atribuído aos conceitos 1 e 5.
Dos 62 cursos da região Sul, 47 eram de Universidades, para os quais o conceito modal foi 4, atribuído a 22 cursos.
Nesse tipo de organização, nenhum curso ficou sem conceito ( SC ) ou recebeu conceito 1.
Os demais cursos receberam os conceitos 2 ( três cursos ), 3 ( 12 cursos ) e 5 ( dez cursos ).
Os Centros Universitários da região Sul contaram com seis cursos, que obtiveram os conceitos 3 ( quatro cursos, conceito modal ), 4 ( um curso ) e 5 ( um curso ).
Nenhum curso ficou sem conceito ( SC ) ou recebeu os conceitos restantes na categoria.
Dos nove cursos oferecidos em Faculdades na região Centro-Oeste, seis receberam o conceito modal 4.
Aos demais cursos foram atribuídos os conceitos 2 ( um curso ) e 3 ( dois cursos ).
Nenhuma curso da Área de Agronomia em Faculdades da região Sul ficou sem conceito ( SC ).
Na região Centro-Oeste, 29 dos 44 cursos eram de Universidades.
Nesse tipo de organização, nenhum curso recebeu conceito 5.
Os conceitos 1 e 2 foram atribuídos a apenas um curso cada e o conceito 3, conceito modal, a 19 cursos.
O conceito 4, por sua vez, foi atribuído a oito cursos.
Dos quatro cursos participantes de Centros Universitários da região Centro-Oeste, nenhum ficou sem conceito ( SC ) ou recebeu conceito 1 ou 5.
O conceito modal, 2, foi atribuído a dois cursos.
Os conceitos restantes, 3 e 4, foram atribuídos, a um curso cada.
Quanto aos 11 cursos vinculados a Faculdades na região Centro-Oeste, sete receberam o conceito modal 2.
O conceito 1 foi atribuído a um curso, assim como o conceito 4, e o conceito 3, a dois cursos.
Nenhum curso recebeu conceito 1, tampouco obteve conceito 5.
O Gráfico 5.3 apresenta a distribuição do Conceito Enade segundo a Organização Acadêmica da IES.
Os cursos em Universidades ( linha azul ) apresentam uma poligonal mais à direita do que as demais, denotando uma melhor distribuição de Conceitos Enade.
No outro extremo, temos os cursos oferecidos em Faculdades ( linha vermelha ) com a poligonal mais à esquerda e, consequentemente, uma distribuição de Conceitos Enade com menores valores.
IES de Centros Universitários ( linha verde ) apresentam uma situação intermediária.

% matched lemmas: 
\end{textsample}
